% !TEX root = ../main.tex
% =====================================================================
%  config/macros.tex — LỆNH TỰ ĐỊNH NGHĨA DÙNG CHUNG
% ---------------------------------------------------------------------
%  * Chỉ NHÓM TRƯỞNG chỉnh sửa file này.
%  * Thành viên KHÔNG dùng \newcommand / \def trong file chương.
%    Cần macro mới -> đề xuất với nhóm trưởng để thêm vào đây.
% =====================================================================

% ---------------------------------------------------------------------
% \chenanh{<độ rộng>}{<đường dẫn ảnh>}  (lệnh nội bộ, thường không gọi trực tiếp)
%   Nếu ảnh chưa có -> hiện khung đỏ "THIẾU ẢNH" thay vì làm lỗi biên dịch.
%   => Đường dẫn PHẢI có phần mở rộng (.png / .jpg / .pdf).
% ---------------------------------------------------------------------
\newcommand{\chenanh}[2]{%
    \IfFileExists{image/#2}%
        {\includegraphics[width=#1]{#2}}%
        {\fcolorbox{red}{red!5}{\parbox[c][3cm][c]{\dimexpr #1-2\fboxsep-2\fboxrule\relax}{\centering\color{red}%
            \textbf{THIẾU ẢNH}\\\footnotesize\texttt{image/\detokenize{#2}}}}}%
}

% ---------------------------------------------------------------------
% \hinh[<độ rộng>]{<đường dẫn ảnh>}{<chú thích>}{<label>}
%   Chèn 1 hình, căn giữa, cố định vị trí [H].
%   - Độ rộng mặc định 0.8\textwidth.
%   - Đường dẫn tính từ thư mục image/  (VD: chuong3/vi_converter.png)
%   VD: \hinh[0.6\textwidth]{chuong3/vi_converter.png}{Mạch chuyển đổi V-I}{fig:c3:vi-converter}
% ---------------------------------------------------------------------
\newcommand{\hinh}[4][0.8\textwidth]{%
    \begin{figure}[H]
        \centering
        \chenanh{#1}{#2}
        \caption{#3}
        \label{#4}
    \end{figure}%
}

% ---------------------------------------------------------------------
% \haihinh{<ảnh 1>}{<chú thích 1>}{<label 1>}{<ảnh 2>}{<chú thích 2>}{<label 2>}
%   Hai hình đặt cạnh nhau, mỗi hình có số hình riêng (Hình x.1, Hình x.2).
% ---------------------------------------------------------------------
\newcommand{\haihinh}[6]{%
    \begin{figure}[H]
        \centering
        \begin{minipage}[t]{0.47\textwidth}
            \centering
            \chenanh{\textwidth}{#1}
            \caption{#2}
            \label{#3}
        \end{minipage}\hfill
        \begin{minipage}[t]{0.47\textwidth}
            \centering
            \chenanh{\textwidth}{#4}
            \caption{#5}
            \label{#6}
        \end{minipage}
    \end{figure}%
}

% ---------------------------------------------------------------------
% \nhanxet{<nội dung>}  — đoạn nhận xét sau mỗi kết quả mô phỏng / đo đạc
% ---------------------------------------------------------------------
\newcommand{\nhanxet}[1]{%
    \par\smallskip\noindent\textbf{Nhận xét:} #1\par\smallskip
}

% ---------------------------------------------------------------------
% \ketqua{<công thức>}  — đóng khung kết quả tính toán quan trọng (dùng trong môi trường toán)
%   VD: \[ \ketqua{R_G = \qty{100}{\ohm}} \]
% ---------------------------------------------------------------------
\newcommand{\ketqua}[1]{\boxed{#1}}

% ---------------------------------------------------------------------
% \tieumuc{<tiêu đề>}  — tiêu đề in đậm không đánh số, không vào mục lục
%   Dùng cho các ý nhỏ hơn subsection (VD: "Cấu tạo của Load cell:")
% ---------------------------------------------------------------------
\newcommand{\tieumuc}[1]{\par\medskip\noindent\textbf{#1}\par\nopagebreak}

% ---------------------------------------------------------------------
% \TODO{<ghi chú>}  — đánh dấu phần chưa hoàn thành (chữ đỏ, dễ thấy khi review)
%   PHẢI xóa hết \TODO trước khi merge vào bản final.
% ---------------------------------------------------------------------
\DeclareRobustCommand{\TODO}[1]{\textcolor{red}{\textbf{[TODO: #1]}}}

% ---------------------------------------------------------------------
% Ký hiệu hay dùng
% ---------------------------------------------------------------------
\newcommand{\Vin}{V_{\mathrm{in}}}
\newcommand{\Vout}{V_{\mathrm{out}}}
\newcommand{\Vref}{V_{\mathrm{ref}}}
\newcommand{\Rshunt}{R_{\mathrm{shunt}}}
